\chapter{Find the boundary of a nested value}
\MARSChapterHeading{Find the boundary of a nested value}
\label{revision:lab-cbor}

A parser can know what a value claims to contain before knowing where that value
ends. A CBOR head may describe a four-pair map in two bytes; the pairs still have
to be traversed. This lab follows that distinction through nested calls, complete
record acceptance and preferred output, then asks you to count top-level items.

Import the \texttt{032-cbor-telemetry} student archive and open it in Learn.
The complete entry \texttt{cbor-telemetry.asm} selects a head read. Its
\texttt{worked/} copies select one item, decoding, or the complete memory
pipeline. The independent sequence entry is deliberately unfinished. Read the
main chapter's head, continuation and projection explanations first.

Answer predictions inline. Their source pointers select the intended worked
copy. During Debug, inspect the selected paused frame: Learn can outline the next
requested boundary after recognizing the current one. A highlighted instruction
has not executed yet. Use one breakpoint at a time, resume between observations,
and finish each process before opening another source. Repeated labels in the
recursive walk refer to different visits and sometimes different active frames.
The observations describe these supplied programs, not a required register
allocation for your independent implementation.

\section{Read a head without inventing its payload}

The selected trace begins \texttt{b8 04}. Split \texttt{b8} into its high three
bits and low five. The major type is five, meaning map; the low field is 24,
meaning that one following byte supplies the argument. That byte is four, so
the head describes four pairs and occupies two bytes. Four is not a message
length, and two is not the complete map extent.

At \texttt{after\_h\_type}, \texttt{\$t1} contains five while the initial
prefix remains in \texttt{\$t0}. At \texttt{head\_return}, inspect status,
extent and argument separately: \texttt{\$v0}=0, \texttt{\$v1}=2,
\texttt{\$a2}=5 and \texttt{\$a3}=4. The caller has not yet stored the latter
two values in \texttt{observed\_major} and \texttt{observed\_argument}.
Those stores happen after this return checkpoint.

The truncated-head copy still physically contains \texttt{b8 04}, but selects
only the first byte. Predict its result before running. The required argument
byte is outside the selected interval. At \texttt{head\_return}, expect status
one and zero in \texttt{v1/a2/a3}. A readable allocation is not permission to
use unselected bytes, and a missing argument does not mean an encoded zero.

The full-counter copy supplies \texttt{1b 01 02 03 04 05 06 07 08} with count
nine. At \texttt{after\_h\_width}, eight argument bytes plus the initial byte
fit. The loop shifts the accumulator eight bits left with \texttt{dsll} and ORs
in an unsigned byte. After four bytes it is \texttt{01020304}; the fifth produces
\texttt{0102030405}, which no longer fits a word. At
\texttt{after\_h\_bytes}, both the accumulator and \texttt{\$a3} retain
\texttt{0102030405060708}, and the head extent is nine.

For an additional diagnosis, single-step that loop in a disposable copy with
\texttt{sll} replacing \texttt{dsll}. Locate the first divergent iteration,
not merely the final low word. Restore the source before continuing. The bytes
have constructed a numeric argument; its later native storage order is a separate
question.

\section{Resume a parent with its own remaining work}

The nested-item copy selects ten bytes:
\begin{verbatim}
a2 01 82 42 80 ff 21 00 80 00
\end{verbatim}
The first nine form a map. Key one has an array value containing a two-byte byte
string and a negative integer. Key zero has an empty-array value. The final zero
is an integer sibling after that map. A walker for one item must return extent
nine, leaving the sibling unread.

Before following the recursive calls, identify what belongs to each suspended
parent. A return address finds the next instruction, but the parent also needs
its cursor, remaining byte budget, outstanding children or pairs, depth and map
key identity. The child cannot supply those values merely by returning its own
extent.

\begin{table}[htbp]
\centering
\begin{tabular}{l r r r}
\toprule
Active parent boundary & Work left & Byte budget & Child depth\\
\midrule
Map after key one & 2 pairs & 8 & 1\\
Array before its first child & 2 children & 7 & 2\\
Array after the byte string & 1 child & 4 & 2\\
Array after the negative integer & 0 children & 3 & 2\\
Map after the complete array & 1 pair & 3 & 1\\
\bottomrule
\end{tabular}
\caption{The supplied recursive walker restores the appropriate parent's
continuation. Byte budgets still include the final sibling in the selection;
they do not mean every remaining byte belongs to this container. Work is decremented
after a complete child or pair value returns.}
\label{tab:revision-cbor-lab-continuations}
\end{table}

Follow \texttt{after\_k\_head}, \texttt{after\_k\_key},
\texttt{after\_k\_array} and \texttt{after\_k\_string}. The string head
\texttt{42} admits two payload bytes, \texttt{80 ff}. They are data, not an
empty array followed by a control value. Consuming that string leaves four
selected bytes from the child's cursor.

The next three requested \texttt{after\_k\_child} visits show why the active
frame matters. First the array resumes with one child left. Next it resumes
with none. Finally the map resumes with one pair left and its key-one mask
still present. The helper's \texttt{\$a2} is major one after the negative child,
while the array's saved child depth in \texttt{\$s2} remains two. In the restored
map, that saved depth is one. The same register has a routine-specific meaning;
an argument register may have acquired a new meaning after a call.

At \texttt{skip\_return}, confirm nine consumed bytes. The remaining integer
belongs to a sequence loop, not to this one-item walk. Counting the map's nested
children as separate top-level items would confuse these two interfaces.

\section{Keep depth admission separate from key identity}

At inherited depth eight, the empty-array copy has a readable, valid head.
At \texttt{after\_k\_head}, its major type is four and its argument is zero.
It nevertheless cannot open another container. At \texttt{skip\_return}, expect
status one and zero extent. Emptiness does not remove the container level.

Use the supplied depth-seven copy for a contrasting manual run: the empty array
is permitted and consumes one byte. Then use the primitive-eight copy. Its
integer consumes one byte without opening a container. The depth rule therefore
cannot reject every value at eight before distinguishing its type. Siblings
also do not accumulate depth from one another; the independent loop must retain
the original inherited depth across helper calls.

The duplicate-key copy selects \texttt{a2 00 00 18 00 00}. Its two key heads
have different widths but both identify numeric zero. At \texttt{skip\_bad},
bit zero is already set in the current map's mask. The last head may still have
status zero before the highlighted \texttt{li} constructs the map failure.
At the return, that failure must have zero extent, not the successful key's width.

Contrast the valid nested-map copy \texttt{a1 00 a1 00 00}. Both the outer and
inner maps use key zero. Observe \texttt{after\_k\_key} first in the outer
frame and then in the inner frame. Each mask contains bit zero, but each belongs
to its own map. The child depths are one and two. A global mask would reject
this valid item; comparing encoded key byte strings would admit the invalid
repeated key. These are different mistakes.

\section{Publish a complete local record}

The complete decode copy selects the original 52-byte trace. At
\texttt{after\_d\_finish}, all pairs and selected bytes have been consumed,
and the seen mask contains keys zero through three. Its low three bits establish
the required fields. The private candidate exists, but the caller record is
still untouched. This is the boundary that authorizes publication.

At \texttt{after\_d\_payload}, the active fields have been copied and the
result is success with extent 52. At \texttt{decode\_return}, inspect the
local record. Its counter is \texttt{0102030405060708}, identifier length is
three and reading count is ten. The identifier bytes are \texttt{00 ff 41}.

\Needspace{18\baselineskip}
\begin{figure}[!ht]
  \centering
  \begin{tikzpicture}[alt={Counter, lengths and binary identifier in native memory. The counter is stored little-endian, lengths are three and ten, and the next row begins 00 ff 41 before unchanged sentinel bytes.}]
    \node[anchor=north west,inner sep=0,font=\small\bfseries] (title) at (0,0) {Counter, lengths and counted identifier};
    \node[anchor=north west,inner sep=0] (a) at ([yshift=-2mm]title.south west) {\MARSRevisionCrop{revision/captures/cbor/cbor-record.png}{0.99\linewidth}{0.03125}{.31555556}{0.58203125}{0.60888889}};
    \begin{scope}[shift={(a.south west)},x={($(a.south east)-(a.south west)$)},y={($(a.north west)-(a.south west)$)}]
      \draw[black,line width=.7pt] (.20,.3088) rectangle (.58,.6176);\draw[black,line width=.7pt] (.20,.005) rectangle (.345,.303);
    \end{scope}
  \end{tikzpicture}
  \caption{The counter occupies eight native little-endian bytes, followed by lengths three and ten. The next row begins with identifier bytes 00 ff 41; its unused tail remains a5. Native LaTeX supplies the title, crop and outlines.}
  \label{fig:revision-cbor-native-record}
\end{figure}


The counter's native bytes are \texttt{08 07 06 05 04 03 02 01}, while its
wire argument was big-endian. The identifier does not undergo a counter-style
reversal: its bytes are binary data, including the first zero, without a
terminating-zero convention.

The ten readings are \(0,23,24,-1,-24,-25,255,256,32767,-32768\).
At record offset 48, inspect signed halfwords or pairs of bytes. Minus one has
local bits \texttt{ffff}; minus 32768 has \texttt{8000}. Those are native
signed codes, not CBOR heads. For negative major type one, argument 32767 means
$-1-32767=-32768$ before storage.

\Needspace{18\baselineskip}
\begin{figure}[!ht]
\centering
\begin{tikzpicture}[alt={Ten decoded signed 16-bit readings across two native Memory rows. An outline marks ff 7f and 00 80, the two range endpoints, followed by unchanged a5 bytes.}]
\node[anchor=north west,inner sep=0,font=\small\bfseries] (title) at (0,0) {Ten signed readings and an untouched tail};
\node[anchor=north west,inner sep=0] (a) at ([yshift=-2mm]title.south west) {\MARSRevisionCrop{revision/captures/cbor/cbor-samples.png}{0.99\linewidth}{0.03125}{.31555556}{0.58203125}{0.60888889}};
\begin{scope}[shift={(a.south west)},x={($(a.south east)-(a.south west)$)},y={($(a.north west)-(a.south west)$)}]
\draw[black,line width=.7pt] (.20,.005) rectangle (.39,.303);
\end{scope}
\end{tikzpicture}
\caption{The final two readings are 32767 and minus 32768, stored as little-endian bytes ff 7f and 00 80. The remaining bytes retain a5. These are local halfword codes rather than CBOR heads. Native LaTeX supplies the title, crop and outlines.}
\label{fig:revision-cbor-native-samples}
\end{figure}


Inspect the unused identifier bytes after offset 18 and unused sample bytes after
67. They retain their \texttt{a5} pattern. Success writes active fields rather
than clearing every allocated byte. The guard bytes beyond the 176-byte record
interface also remain unchanged.

The late-counter copy offers only 51 bytes. Its final counter head lacks one
argument byte, although that byte is still physically present. Earlier fields
can look convincing, but the whole caller record remains unchanged. At
\texttt{decode\_return}, expect status one and zero extent.

\begin{figure}[!ht]
\centering
\begin{tikzpicture}[alt={Native record Memory after the telemetry decoder rejects a counter missing its final argument byte. The displayed rows remain a5 sentinels; the complete caller record is checked separately.}]
\node[anchor=north west,inner sep=0,font=\small\bfseries] (title) at (0,0) {A late missing byte prevents record publication};
\node[anchor=north west,inner sep=0] (a) at ([yshift=-2mm]title.south west) {\MARSRevisionCrop{revision/captures/cbor/cbor-late.png}{0.99\linewidth}{0.03125}{.31555556}{0.58203125}{0.60888889}};
\begin{scope}[shift={(a.south west)},x={($(a.south east)-(a.south west)$)},y={($(a.north west)-(a.south west)$)}]
\draw[black,line width=.7pt] (.20,.005) rectangle (.965,.6176);
\end{scope}
\end{tikzpicture}
\caption{The final counter argument is truncated at selected count 51. Both displayed record rows retain their initial a5 bytes. Separate access and whole-arena checks establish preservation beyond the pictured prefix. Native LaTeX supplies the title, crop and outlines.}
\label{fig:revision-cbor-native-late}
\end{figure}


The pictured prefix helps locate the record; executable checks separately inspect
the complete arena. As another experiment, append integer zero after the complete
root and select 53 bytes. The walker can identify a valid 52-byte first item,
but complete telemetry decoding must reject the extra item. Valid generic syntax
and this application's complete-record contract are separate requirements.

\section{Derive the preferred output and its publication boundary}

The pipeline copy selects head, walk, decode, exact bound and encode in order.
It obtains valid record input through decoding; switching directly to an
encode-only caller would not prepare that input automatically. Derive the exact
37-byte output before stopping at \texttt{bound\_return}:
\[
 4+9+1+3+1+19=37.
\]
The terms are fixed map/key bytes, the counter head and argument, identifier head,
identifier payload, array head, and reading encodings. This measures this record,
not the maximum 241-byte reservation for every local record.

The input used 52 bytes. The output drops the validated optional field and chooses
preferred heads for the required fields. It starts
\texttt{a3 00 1b 01 02 03 04 05 06 07 08}, followed by
\texttt{01 43 00 ff 41} and the readings. Input reading 23 used
\texttt{18 17}; output uses \texttt{17}. Reading 24 still needs
\texttt{18 18}. Explain why both follow the same shortest-head rule.

\begin{figure}[!ht]
\centering
\begin{tikzpicture}[alt={Native encoded output, exactly 37 bytes across three sixteen-byte rows. The third row contains five encoded bytes followed by unchanged a5 storage. Native outlines distinguish the selected byte sequence from its tail.}]
\node[anchor=north west,inner sep=0,font=\small\bfseries] (title) at (0,0) {Thirty-seven bytes, then unchanged storage};
\node[anchor=north west,inner sep=0] (a) at ([yshift=-2mm]title.south west) {\MARSRevisionCrop{revision/captures/cbor/cbor-output.png}{0.99\linewidth}{0.03125}{.29222222}{0.58203125}{0.60888889}};
\begin{scope}[shift={(a.south west)},x={($(a.south east)-(a.south west)$)},y={($(a.north west)-(a.south west)$)}]
\draw[black,line width=.7pt] (.20,.012) rectangle (.44,.235);
\end{scope}
\end{tikzpicture}
\caption{The output begins with the three-pair map and key-zero counter. Its third row ends after five encoded bytes; the next a5 byte is outside the returned count. Optional metadata and wider input spellings are absent from this projection. Native LaTeX supplies the title, crop and outlines.}
\label{fig:revision-cbor-native-output}
\end{figure}


At \texttt{encode\_return}, expect success and 37 bytes. At
\texttt{results\_ready}, those values are still live before the caller stores
them. The output preserves required field meanings, but it is not a byte-for-byte
round trip of the input map or a lossless encoding of its optional metadata.
Generic-decoder interoperability checks the intended projection separately.

The short-capacity copy supplies 36 output bytes. At \texttt{decode\_return},
the local record has already been published successfully. At
\texttt{after\_e\_bound}, the needed size is 37 and the saved capacity is 36;
the destination is still untouched. At \texttt{encode\_return}, expect capacity
status two and zero length. Decode's record remains valid. This caller does not
promise whole-pipeline rollback, while encoding still promises no partial output
on its own capacity failure.

The silent ordinary Run uses the head entry. Its normal exit does not establish
a complete accepted telemetry record by itself. Retain the numerical Debug
observations and memory contents at their intended boundaries. Additional complete
codec checks remain available in the public book repository:
\begin{verbatim}
java -jar /path/to/cli.jar grade \
  packages/032-cbor-telemetry/checks/decode.kts \
  revision/labs/files/cbor-telemetry.asm
\end{verbatim}
Use head, skip, format, bound, encode and caller checks for their corresponding
worked interfaces. These are finite checks, not a proof of every nested structure.

\section{Count top-level items in your own scanner}

Complete only \texttt{cbor\_sequence\_count} in the delivered
\texttt{independent/cbor-sequence.asm}. Keep its caller and complete helpers.
Its default mode eight supplies the trace as one top-level item. Read the exact
full-width interface and early-error order in \texttt{independent/README.md}.

Begin with \texttt{00 82 01 02 40}. Draw three intervals: an integer, an array
including both children, and an empty byte string. Derive count three, then
identify what survives each walker call. Retain the original inherited depth for
every sibling. Advance by a positive returned byte extent, not by a head width,
and count one item per complete return. Prove that each advance fits and reduces
the remaining selected interval.

A valid empty selection succeeds without using an unusable pointer. Excessive
unsigned count fails before invalid inherited depth, and neither needs a read.
Try a valid first item followed by \texttt{1b}: the incomplete suffix must clear
the aggregate result, rather than retaining credit for the valid prefix. A later
empty call must also establish its own zero count instead of reusing a previous
successful result register.

The seven native groups preserve the original 203 cases and add real caller
termination and a same-session retry. Their hints distinguish interval arithmetic,
saved continuation, failure result and depth errors. The retry keeps one session
through rejection, success, empty input and a sibling depth failure. The caller
check verifies the stored result after return and actual exit. The untouched body
fails deliberately; it is not a simulator defect.

The optional CLI route retains the original task:
\begin{verbatim}
java -jar /path/to/cli.jar grade \
  revision/labs/checks/cbor-sequence.kts \
  /path/to/your/cbor-sequence.asm
\end{verbatim}
Equivalent solutions can retain remaining bytes or an end address, reuse the
walker or traverse independently, and choose another frame layout. Keep a parent
continuation trace, a source-permission argument and counterexamples for sibling
depth, separate map identities, malformed suffixes and selected boundaries. Those
make the result explainable when the supplied trace is no longer the test input.

\section{Companion lab: evaluate an expression language}
\label{revision-lab:expression}

CBOR makes nested structure explicit in its heads and container counts. A text
expression instead uses punctuation, operator rules and a byte cursor. This lab
asks you to implement a small language whose result depends on those rules, not
merely to replace characters with arithmetic instructions. It revises original
project \texttt{017-expression-evaluator}.

Use \texttt{expression-evaluator.asm} in \texttt{revision/labs/files} as your
working skeleton. Implement \texttt{eval\_expr}; the entry deliberately stops with
\texttt{break} until replaced. Keep the supplied caller and data arenas. The
separate complete \texttt{expression-arithmetic.asm} demonstrates checked
operations without supplying the parser. Original grammar and feedback remain
in \texttt{packages/017-expression-evaluator}.

The prerequisites are bounded byte access, signed-word representation, full-width
function frames and the parent-state reasoning in the nested-data chapter.
The new work connects grammar to control flow, retains pending operations across
calls, and reports the first failed requirement precisely. These activities use
ordinary Run and Debug; their native guided lesson remains unfinished.

\subsection*{Give the input a grammar before writing the loop}

For \texttt{18 - 3 * (2 + -4)}, first evaluate the group to minus two, then the
product to minus six, then subtract that product from eighteen. The result is
24. A loop that immediately applies each operator to the next number would
instead reach fifteen before seeing multiplication. Operator precedence tells
you which larger operand must complete before an operation can run.

\begin{figure}[htbp]
  \centering
  \begin{minipage}{.94\linewidth}
    \begin{verbatim}
expression := term { ('+' | '-') term }
term       := unary { ('*' | '/') unary }
unary      := ('+' | '-') unary | primary
primary    := digits | '(' expression ')'
digits     := one or more ASCII decimal digits
\end{verbatim}
  \end{minipage}
  \caption{The exercise language. Braces mean zero or more repetitions; the vertical bar separates alternatives. Quoted characters are actual input bytes. Parentheses outside quotes group the grammar notation. Whitespace may occur between tokens, but cannot join two digit sequences.}
  \label{fig:revision-expression-grammar}
\end{figure}


Read the expression rule as: parse one term, then repeatedly accept a plus or
minus followed by another complete term. The term rule does the same for
multiplication and division, using unary values as operands. A unary sign owns
a whole following unary value, so \texttt{--5} means the negation of minus five.
A primary is either a decimal literal or a group containing another expression.
The group owns both parentheses; the inner expression does not consume its
caller's closing parenthesis.

The loops implement left association. In \texttt{9-6/2-1}, the first right term
is \texttt{6/2}, giving three; the expression loop computes \(9-3=6\), then
\(6-1=5\). A recursive call for an unrestricted expression after each minus
could consume the rest as its right operand and change that association.
One routine per precedence level is a useful recursive-descent design, not a
required private organization.\footnote{For background on grammar-driven
recursive descent, see Robert Nystrom, \emph{Crafting Interpreters},
\href{https://craftinginterpreters.com/parsing-expressions.html}{Parsing expressions}.
This exercise defines its own grammar, numeric bounds and diagnostic order.}

Accept only ASCII digits, the six operator/grouping characters, and space,
tab, carriage return or line feed between tokens. The input is a selected byte
slice, not a zero-terminated string. An embedded zero is invalid. A vertical tab
is also invalid, even if another language calls it whitespace. \texttt{12 34}
contains two literals without an intervening operator; whitespace does not
turn it into 1234. No identifiers, comments, implicit multiplication, hexadecimal
or floating-point forms belong to this language.

\subsection*{Make the cursor and suspended work explicit}

Represent the input by a full base address, current byte offset and authoritative
length. Before any load, establish that the complete load lies in the slice.
An unsigned byte load is enough to distinguish ASCII punctuation and digits.
Reaching the end while skipping whitespace is ordinary progress; reaching it
while waiting for a primary or closing parenthesis is a syntax failure.

A useful private helper result is a value together with an advanced shared
cursor and an error state. When a term calls a unary parser, it must retain its
left value, pending operator and that operator's source offset. The recursive
child advances the cursor. Restoring the old cursor on return would undo the
child's progress and can repeat the same token indefinitely.

\begin{table}[htbp]
  \centering
  \begin{tabular}{r l l}
    \toprule
    Cursor offset & Completed work                & Pending operation             \\
    \midrule
    2             & Literal 18                    & Outer subtraction begins at 3 \\
    6             & Literal 3                     & Multiplication begins at 7    \\
    11            & Group's left term 2           & Addition begins at 12         \\
    16            & Unary minus has produced $-4$ & Group computes $2+(-4)$       \\
    17            & Group closes with value $-2$  & Term computes $3(-2)$         \\
    17            & Term returns $-6$             & Expression computes $18-(-6)$ \\
    \bottomrule
  \end{tabular}
  \caption{A trace of \texttt{18 - 3 * (2 + -4)}. Offsets are zero-based byte positions; 17 is the input end. A returned value and the cursor are different outputs of parsing. Whitespace skipping can advance the cursor between these boundaries.}
  \label{tab:revision-expression-trace}
\end{table}


The term awaiting the group in the worked input retains left value three,
operator \texttt{*} and offset seven. The group creates a nested expression
whose pending addition retains two and offset twelve. The unary parser remembers
the minus sign at fourteen while reading literal four. After that child
returns, negation gives minus four; the group gives minus two; the waiting
multiplication can finally run. Each frame owns its pending operation while
the cursor records shared forward progress.

Distinguish the public calling convention from private parser state. The public
\texttt{eval\_expr} must preserve every bit of \texttt{s0} through \texttt{s7},
\texttt{fp} and \texttt{sp}, and return through its incoming \texttt{ra}. A private
helper convention may deliberately share some saved registers as parser state,
provided the public wrapper restores the original caller values. Document that
convention; mechanically restoring every shared parse register in every private
helper would erase progress or the first error. Save return addresses and live
operands in full-width, balanced, properly aligned frames.

Another valid design uses a precedence threshold and one binary-expression
routine. After an operator of precedence $p$, parse its right operand with a
threshold excluding another operator of that same precedence; the enclosing
loop then performs equal-precedence operations left to right. This is often
called precedence climbing. Choose a design you can trace. Equivalent private
helpers and explicit stacks are accepted when they preserve the whole contract.

\subsection*{Compute widely, then prove the signed-word range}

Decimal literals are limited to zero through 2147483647. For current literal
value $n$ and next digit $d$, form $10n+d$ in a doubleword and test the result
before accepting that digit. The previous accepted value was at most 2147483647,
so this one-digit candidate fits comfortably in 64 bits. Record the current
digit's offset if it exceeds the positive limit. Long leading zeros are legal
within the 64-byte input bound.

Every intermediate arithmetic value must fit signed 32 bits, from minus
2147483648 through 2147483647. A wide candidate prevents word arithmetic from
hiding overflow before the test. Addition and subtraction of two word values
fit a signed doubleword; their product has magnitude at most \(2^{62}\).
After excluding zero division, a wide quotient also fits. The special quotient
\(-2147483648/-1=2147483648\) is a valid wide result that must fail the narrower
language range.

\Needspace{17\baselineskip}
Run the complete arithmetic example and break at \texttt{product\_ready}, then
\texttt{arithmetic\_ready}. Its first call computes \(3(-2)=-6\); the second
computes \(18-(-6)=24\). Both return zero status. At
\texttt{arithmetic\_stored}, the three native words hold product minus six,
status zero and result 24. The caller does not print them. Both supplied callers round the initial stack
pointer down to a sixteen-byte boundary before their first call.

\begin{figure}[!ht]
  \centering
  \begin{tikzpicture}[alt={Native Memory observation of the complete public checked-arithmetic example. An outline marks three 32-bit words: product minus six, zero status and result 24. These correspond to three times minus two, followed by eighteen minus that product.}]
    \node[anchor=north west,inner sep=0,font=\small\bfseries] (title) at (0,0) {A checked product feeds the next operation};
    \node[anchor=north west,inner sep=0] (a) at ([yshift=-2mm]title.south west) {\MARSRevisionCrop{revision/captures/expression/expression-arithmetic.png}{0.99\linewidth}{0.03125}{0.33888889}{0.58203125}{0.60888889}};
    \begin{scope}[shift={(a.south west)},x={($(a.south east)-(a.south west)$)},y={($(a.north west)-(a.south west)$)}]
      \draw[black,line width=.7pt] (.20,.005) rectangle (.77,.447);
    \end{scope}
  \end{tikzpicture}
  \caption{The first three little-endian words hold product minus six, status zero and final result 24. They are stored at \texttt{arithmetic\_stored} in the complete public arithmetic example. Native LaTeX supplies the title, crop and outline.}
  \label{fig:revision-expression-native-arithmetic}
\end{figure}


\Needspace{14\baselineskip}
In \texttt{checked\_binary}, \texttt{a0/a1} contain signed-word operands,
\texttt{a2} is the ASCII code for one of \texttt{+-*/}, and \texttt{a3} is that
operator's byte offset. It returns status in \texttt{v0} and value or error
offset in \texttt{v1}. The complete helper checks zero before executing
\texttt{ddiv}. Once \texttt{t0} holds a wide arithmetic candidate, its range
comparison begins:
\Needspace{8\baselineskip}
\MARSCodeLabel{Compare a wide candidate with its sign-extended low word}
\begin{minted}{mars-idea}
    sll $t1, $t0, 0
    beqc $t1, $t0, checked_value
    nop
\end{minted}

Here the word-sized \texttt{sll} intentionally discards the upper half and
sign-extends the low word. Equality with the original doubleword proves that
no information needed for a signed-word value was lost. It is a test using a
copy, not permission to replace the original candidate with a wrapped word.
For 2147483648, the copy is minus 2147483648, so equality fails.

Use \texttt{revision/labs/checks/expression-arithmetic.kts} to check the worked
helper and caller. Its operands and operators are already supplied as numbers;
it neither parses a token nor diagnoses malformed text. Your evaluator may
reuse or adapt it, with an explicit private interface, or implement equivalent
checked arithmetic differently.

\subsection*{Treat diagnostics as part of evaluation}

Return status zero and the signed value on success. Failure status minus one
means syntax, minus two overflow, and minus three division by zero. On failure,
\texttt{v1} changes meaning: it is a zero-based byte offset. The offset can equal
the input length when an item is missing at the end. Do not confuse a successful
numeric result such as ten with an error offset ten; inspect status first.

\begin{table}[htbp]
  \centering
  \begin{tabular}{l r r l}
    \toprule
    Selected input            & Status & Offset & First failed requirement       \\
    \midrule
    \texttt{2147483647+1?}    & $-2$   & 10     & Addition's word range          \\
    \texttt{2147483647+1*?}   & $-1$   & 13     & Multiplication's right operand \\
    \texttt{1/(2-2)}          & $-3$   & 1      & Completed denominator is zero  \\
    \texttt{1/(2-2}           & $-1$   & 6      & Group needs its closing byte   \\
    \texttt{-2147483648}      & $-2$   & 10     & Literal exceeds positive bound \\
    \texttt{-(-2147483647-1)} & $-2$   & 0      & Outer negation's word range    \\
    \bottomrule
  \end{tabular}
  \caption{Error identity follows the specified parse-and-evaluate procedure. A later invalid byte does not replace an earlier arithmetic failure; an operation cannot fail before its required operand has successfully completed.}
  \label{tab:revision-expression-errors}
\end{table}


In \texttt{2147483647+1?}, the right term one has completed, and the addition
fails at offset ten before final consumption checks the question mark. With
\texttt{2147483647+1*?}, that right term cannot finish: multiplication still
needs its right unary operand at thirteen. The syntax failure therefore occurs
before the outer addition can be attempted. A whole-input syntax prevalidation
pass would change the required diagnosis in the first example.

The following diagnostic figures use a completed evaluator with the same public
caller and data layout as your skeleton. At \texttt{expression\_stored}, the
row begins with length, status and value-or-offset, followed by a guard.

\begin{figure}[!ht]
\centering
\begin{tikzpicture}[alt={Actual native Memory after a completed evaluator processes 2147483647+1?. The row holds length thirteen, overflow status minus two, byte offset ten and an unchanged guard. An outline marks status and offset.}]
\node[anchor=north west,inner sep=0,font=\small\bfseries] (title) at (0,0) {Overflow precedes the trailing invalid byte};
\node[anchor=north west,inner sep=0] (a) at ([yshift=-2mm]title.south west) {\MARSRevisionCrop{revision/captures/expression/expression-overflow.png}{0.99\linewidth}{0.03125}{0.33888889}{0.58203125}{0.60888889}};
\begin{scope}[shift={(a.south west)},x={($(a.south east)-(a.south west)$)},y={($(a.north west)-(a.south west)$)}]
\draw[black,line width=.7pt] (.39,.005) rectangle (.77,.447);
\end{scope}
\end{tikzpicture}
\caption{The caller stores selected length thirteen, status minus two and operator offset ten. The outlined status/offset pair identifies the addition overflow in 2147483647+1?, before complete-consumption checking reaches the question mark. Native LaTeX supplies the title, crop and outline.}
\label{fig:revision-expression-native-overflow}
\end{figure}


Likewise, \texttt{1/(2-2} does not provide a denominator. The inner subtraction
has produced zero, but the group still needs its closing parenthesis at offset
six. Division cannot inspect that provisional group value as a completed operand.
After the closing parenthesis is present, the same denominator completes and
zero division is reported at the slash, offset one.

\begin{figure}[!ht]
  \centering
  \begin{tikzpicture}[alt={Actual native Memory after a completed evaluator processes the six bytes 1/(2-2. The row holds length six, syntax status minus one, byte offset six and an unchanged guard. An outline marks status and offset, distinguishing this from zero division.}]
    \node[anchor=north west,inner sep=0,font=\small\bfseries] (title) at (0,0) {An incomplete group is not a denominator};
    \node[anchor=north west,inner sep=0] (a) at ([yshift=-2mm]title.south west) {\MARSRevisionCrop{revision/captures/expression/expression-missing-close.png}{0.99\linewidth}{0.03125}{0.33888889}{0.58203125}{0.60888889}};
    \begin{scope}[shift={(a.south west)},x={($(a.south east)-(a.south west)$)},y={($(a.north west)-(a.south west)$)}]
      \draw[black,line width=.7pt] (.39,.005) rectangle (.77,.447);
    \end{scope}
  \end{tikzpicture}
  \caption{For the six-byte input 1/(2-2, the stored status is minus one and the offset is six, the missing closing parenthesis at the input end. The division has no completed denominator yet. Native LaTeX supplies the title, crop and outline.}
  \label{fig:revision-expression-native-missing-close}
\end{figure}


The leading minus in \texttt{-2147483648} is a unary operator, not part of the
literal. Its operand literal exceeds 2147483647 at the final digit, offset ten,
so the evaluator never reaches negation. Use \texttt{(-2147483647-1)} to obtain
the minimum through valid literals and an in-range subtraction. Negating that
completed value fails at the outer sign. An eventual multiplication by zero
does not excuse an earlier overflow in \texttt{(2147483647+1)*0}.

Keep the first error while returning through parents. A child failure prevents
its parent's pending arithmetic or closing-byte check from replacing that error.
Only a successful outer expression proceeds to final whitespace skipping and
complete-consumption checking. Useful offset records belong to different events:
a literal remembers the current digit; unary remembers its sign; binary remembers
its operator; missing syntax uses the current cursor or the input end.

\subsection*{Build and investigate your complete evaluator}

The public entry receives a full input pointer in \texttt{a0} and a length from
zero through 64 in \texttt{a1}. The slice lies in a 96-byte arena and may be
unaligned. All caller data, including input and guards, are read-only even for
temporary stores. Read only the selected slice among caller objects. Your own
constant tables and private stack may be used. Produce no console output.
The preserved caller source contains the worked 17-byte expression; it stores
the returned status and result after \texttt{expression\_ready}.

Implement and test boundaries that you can explain. A whitespace skip must
terminate at the slice end. A literal must stop before a nondigit and identify
the first excessive digit. A group must retain its value while checking its
closing byte. A unary sign must wait for a complete operand. A term must finish
multiplication and division before an expression loop uses its value. Connect
those pieces under the public saved-state wrapper, then require complete input
consumption. Intermediate helper tests support construction; they are not final
acceptance of the language.

After your evaluator is complete, break at \texttt{expression\_ready} for the
worked input and expect \texttt{v0}=0, \texttt{v1}=24. At
\texttt{expression\_stored}, inspect its stored words before exit overwrites the
service selector. The native figures use those shared caller boundaries in a
completed implementation; your private helpers need not match its layout.
Change only the selected input and length for each diagnostic case, using a
fresh copy or resetting the machine between ordinary Debug runs.

\Needspace{7\baselineskip}
Run the original public kernel and caller checks against your working source:
\begin{verbatim}
java -jar /path/to/cli.jar grade \
  packages/017-expression-evaluator/checks/kernel.kts \
  /path/to/your/expression-evaluator.asm
\end{verbatim}
Repeat with \texttt{program.kts}. The unfinished skeleton is expected to fail.
The kernel accepts different parsing strategies; it does not prescribe helper
names, frame sizes, traversal order or instruction counts. The caller checks
add observation and stored-result behavior. Passing finite cases does not prove
all 64-byte inputs.

Retain a prediction and explanation for \texttt{9-6/2-1}, consecutive unary
signs, separated digit sequences, the two question-mark examples and the two
denominator groups. Then call the function again after an error using another
input arena. A successful previous call must not seed the next parse, and a
previous failure must not remain as its error state. Compare return status,
value or byte offset, source preservation and balanced caller state; a program
that obtains the right arithmetic value while reading a guard is still incorrect.


\Needspace{13\baselineskip}
\section{Companion lab: interpret a bytecode program}
\label{revision-lab:bytecode}

The expression evaluator reads characters whose grammar determines the next
operation. A bytecode interpreter reads an already encoded sequence of operations.
The bytes still need validation, and their effects belong to a second machine
that your MIPS program maintains. This companion revises original project
\texttt{018-tiny-bytecode-interpreter}.

Import \texttt{018-bytecode-machine}, version 0.1.0, through the MARS Lessons
window. Begin with the complete \texttt{worked/record.asm}, then compare
\texttt{worked/truncated.asm}. Implement \texttt{run\_vm} in
\texttt{independent/vm.asm}; replace its single \texttt{break}, keeping the supplied
caller and data arenas. Private labels, frames and dispatch design are yours to
choose. The original project remains under
\texttt{packages/018-tiny-bytecode-interpreter}.

The preceding nested-data chapter supplies bounds and validation reasoning; the
expression companion supplies checked-word arithmetic. Review functions and stack
preservation as needed. Inline predictions connect the written reasoning to
source pointers. Guided Debug first inspects the complete record classifier,
then the independent VM after you have implemented it. You may also work freely
and return to the staged assessments.

\subsection*{Keep the guest machine distinct from the executing MIPS program}

The MIPS program counter selects instructions such as \texttt{lbu},
\texttt{beqzc} and \texttt{sw}. A \emph{guest PC} is an integer byte offset into
the supplied bytecode image. A guest opcode chooses which native handler to run.
For a guest branch to offset four, your interpreter updates that integer and
returns to its dispatch loop. It does not jump the native machine into the
program's data bytes.

The guest operand stack is another caller-owned data array. The native stack
addressed by \texttt{sp} stores your saved registers and private work. A guest
DUP adds a value to the operand array; it does not duplicate a native return
address. Keep guest depth and native frame size separate even if both happen
to be small integers.

\begin{table}[htbp]
  \centering
  \begin{tabular}{r l l}
    \toprule
    Byte offset & Field                              & Ownership                  \\
    \midrule
    0           & Guest PC                           & Output word                \\
    4           & Live stack depth                   & Output word                \\
    8           & Retired guest count                & Output word                \\
    12          & Status                             & Output word                \\
    16          & Capacity, zero through eight words & Preserved input word       \\
    20          & Fault offset, or $-1$              & Output word                \\
    24          & Full guest-stack address           & Preserved eight-byte input \\
    \bottomrule
  \end{tabular}
  \caption{The eight-byte-aligned descriptor occupies 32 bytes. Output fields report the new invocation; they do not request resumption. The word-aligned guest stack is a separate caller-owned subview. A full stack pointer does not fit in one word.}
  \label{tab:revision-bytecode-descriptor}
\end{table}


The public arguments are full code address, selected length, full descriptor
address and guest instruction budget in \texttt{a0} through \texttt{a3}.
Length and budget each lie between zero and 64. Code is a byte slice inside a
96-byte arena and may be unaligned. The descriptor's input pointer names a
word-aligned stack subview with capacity zero through eight words inside a
sixteen-word allocation. Only the selected capacity is available. Extra
physical cells are not permission to push, read ahead or clear a larger array.

Every call starts at guest PC zero, depth zero and retired count zero. Ignore
the old output fields, even if they describe a previous budget return. The
caller may replace the image at the same address and length. This interface
reports a bounded execution; it does not resume that report on the next call.

\subsection*{Decode instruction boundaries before accepting branches}

\begin{table}[htbp]
\centering
\begin{tabular}{l r l}
\toprule
Opcode and name & Bytes & Guest action\\
\midrule
\texttt{00 HALT} & 1 & Advance, retire and halt\\
\texttt{01 PUSH8} & 2 & Push one signed-byte immediate\\
\texttt{02 PUSH32} & 5 & Push one little-endian signed-word immediate\\
\texttt{10 ADD}, \texttt{11 SUB} & 1 & Replace older left and top right by result\\
\texttt{12 MUL}, \texttt{13 DIV} & 1 & Replace older left and top right by result\\
\texttt{20 DUP}, \texttt{21 DROP} & 1 & Copy or remove top\\
\texttt{22 SWAP}, \texttt{23 OVER} & 1 & Exchange top pair or copy second from top\\
\texttt{30 JMP}, \texttt{31 JZ} & 2 & Absolute byte target; JZ always pops condition\\
\bottomrule
\end{tabular}
\caption{Opcodes are hexadecimal; byte sizes are decimal. Arithmetic uses checked signed words. Branch operands are unsigned byte offsets from the code base. They are neither native instruction addresses nor relative displacements.}
\label{tab:revision-bytecode-encoding}
\end{table}


Walk the entire image from offset zero. At each established instruction start,
read the opcode, determine its size and prove that the complete record fits
before advancing. PUSH8's immediate is a signed byte; PUSH32 carries four
little-endian bytes representing any signed-word bit pattern. Unlike the
expression language's positive decimal literal, PUSH32 can encode minus
2147483648 directly as \texttt{02 00 00 00 80}.

The complete \texttt{worked/record.asm} isolates that first decision. Its
\texttt{record\_width} receives a full pointer, selected length and candidate byte
offset in \texttt{a0/a1/a2}. It proves that the opcode byte exists, classifies its
width, then compares the required width with the remaining selected bytes. On
success it returns status zero and width; on failure it returns a local status
and zero width. This helper does not decode an operand, mark a whole image,
validate a branch target or execute a guest instruction.

Debug the five-byte PUSH32 example. Before \texttt{width\_call}, expect selected
length five and candidate offset zero. At \texttt{record\_fits}, before
\texttt{bltu}, \texttt{t1} holds opcode two, \texttt{t2} holds width five and
\texttt{t3} holds remaining length five. At \texttt{width\_ready}, expect
\texttt{v0}=0 and \texttt{v1}=5. Stop Debug, then open
\texttt{worked/truncated.asm}. Its initializer still contains five physical
bytes, but the selected length is four. Follow the opcode read, then the same
width comparison: required five exceeds remaining four. The returned pair is
\texttt{(-2,0)}. Physical storage past the selected end does not complete the record.

Run the ordinary complete example afterward. It stores status and width and
exits silently. Use \texttt{worked/record.kts} to check all opcode values, each at
selected lengths zero through five, plus the complete caller. These local checks support your
first structural pass; they are not acceptance of the VM.

Record every instruction start. With at most 64 code bytes, a doubleword bitmap
can use bit $p$ for start offset $p$. A private array of 64 byte markers is
another choice. A bitmap needs a doubleword shift, including at offset 63;
a word shift would alias upper-half positions. Do not infer a start merely
because an operand byte happens to equal an opcode.

Only after this structural pass succeeds, scan all encoded branches in image
order. Their unsigned target byte must be below the selected length and belong
to the decoded start set. A target equal to length is invalid. Check unreachable
branches and untaken conditional targets too. A valid runtime path does not
repair an invalid encoded program.

Structural failure takes priority over every bad target. For the bytes
\texttt{30 01 ff}, the first record is a JMP with a bad target into its operand,
but the structural pass next finds invalid opcode \texttt{ff} at offset two.
Report invalid opcode there; target checking never begins. Similarly,
\texttt{00 ff} fails validation after its HALT byte. No guest instruction has
executed, and no guest budget has been consumed in either case.

\Needspace{25\baselineskip}
\begin{table}[!ht]
\centering
\begin{tabular}{r l l}
\toprule
Byte offset & Encoded bytes & Instruction\\
\midrule
0 & \texttt{01 00} & PUSH8 0\\
2 & \texttt{01 03} & PUSH8 3\\
4 & \texttt{20} & DUP\\
5 & \texttt{31 10} & JZ 16\\
7 & \texttt{22} & SWAP\\
8 & \texttt{23} & OVER\\
9 & \texttt{10} & ADD\\
10 & \texttt{22} & SWAP\\
11 & \texttt{01 01} & PUSH8 1\\
13 & \texttt{11} & SUB\\
14 & \texttt{30 04} & JMP 4\\
16 & \texttt{21} & DROP\\
17 & \texttt{00} & HALT\\
\bottomrule
\end{tabular}
\caption{The 18-byte countdown image sums three, two and one. The instruction at offset five has target byte hexadecimal \texttt{10}, decimal 16. Offset six is its operand, even though its byte value is also the ADD opcode. Instruction identity comes from decoding at an established start.}
\label{tab:revision-bytecode-countdown}
\end{table}


Change the target operand at offset six of this image from \texttt{10} to
\texttt{06}. Six is inside the image, but not an instruction start. Even at
zero guest budget, validation must return bad-target status minus three, fault
offset five, PC/depth/retired zero. The fault identifies the branch instruction,
not the operand's location or its target value. This distinction is a useful
check against a range-only validator.

\begin{figure}[!ht]
\centering
\begin{tikzpicture}[alt={Actual native descriptor after zero-budget whole-image validation rejects a branch into an operand. The first three output words are zero, status is minus three and fault is five. Native outlines identify the zero execution prefix and the fault field.}]
\node[anchor=north west,inner sep=0,font=\small\bfseries] (title) at (0,0) {Zero budget still requires a valid target};
\node[anchor=north west,inner sep=0] (a) at ([yshift=-2mm]title.south west) {\MARSRevisionCrop{revision/captures/bytecode/vm-bad-target.png}{0.99\linewidth}{0.03125}{0.31555556}{0.58203125}{0.60888889}};
\begin{scope}[shift={(a.south west)},x={($(a.south east)-(a.south west)$)},y={($(a.north west)-(a.south west)$)}]
\draw[black,line width=.7pt] (.20,.3088) rectangle (.77,.6176);\draw[black,line width=.7pt] (.39,.005) rectangle (.58,.303);
\end{scope}
\end{tikzpicture}
\caption{The image targets operand byte six from the JZ at offset five. PC, depth and retired count remain zero; status is minus three and fault five. The fault names the invalid branch instruction, not its target. No guest instruction has retired. Native LaTeX supplies the title, crop and outlines.}
\label{fig:revision-bytecode-native-bad-target}
\end{figure}


\subsection*{Trace a guest loop through stack effects}

The countdown starts by pushing accumulator zero and counter three. Write the
live stack oldest to newest: \([0,3]\). At offset four, DUP copies the counter,
giving \([0,3,3]\). JZ then consumes the copied condition on either path.
For a nonzero counter, execution falls through to seven with \([0,3]\).
The original counter is still present because DUP made the condition copy.

SWAP gives \([3,0]\), and OVER copies the second value from the top, giving
\([3,0,3]\). ADD consumes the older left operand zero and top right operand
three, producing \([3,3]\). After SWAP, PUSH8 1 and SUB, the stack is
\([3,2]\): accumulated sum three and next counter two. JMP selects offset four.
The next iterations return there with \([5,1]\), then \([6,0]\).

The next DUP/JZ pair takes the exit, leaving \([6,0]\). DROP removes the old
counter, and HALT finishes with the single live value six. A native Memory row
can display inactive words after that live value. Their contents are not more
guest results: the descriptor's depth determines which prefix is meaningful.
The valid countdown needs room for three simultaneous values even though its
final result has depth one.

\begin{figure}[!ht]
\centering
\begin{tikzpicture}[alt={Actual native descriptor after countdown halt. The first row is PC eighteen, depth one, retired thirty-three and status one. The second row retains capacity eight and the full stack pointer and reports fault minus one. Native outline identifies the four main output words.}]
\node[anchor=north west,inner sep=0,font=\small\bfseries] (title) at (0,0) {Halt follows thirty-three guest retirements};
\node[anchor=north west,inner sep=0] (a) at ([yshift=-2mm]title.south west) {\MARSRevisionCrop{revision/captures/bytecode/vm-halted.png}{0.99\linewidth}{0.03125}{0.31555556}{0.58203125}{0.60888889}};
\begin{scope}[shift={(a.south west)},x={($(a.south east)-(a.south west)$)},y={($(a.north west)-(a.south west)$)}]
\draw[black,line width=.7pt] (.20,.3088) rectangle (.965,.6176);
\end{scope}
\end{tikzpicture}
\caption{The first row reports PC 18, depth one, retired count 33 and halt status one. Fault is minus one in the second row. The single live result six is verified separately; the stack cells are not pictured here. Native LaTeX supplies the title, crop and outlines.}
\label{fig:revision-bytecode-native-halted}
\end{figure}


\begin{table}[htbp]
\centering
\begin{tabular}{r r r l l}
\toprule
Guest budget & PC & Retired & Live stack, oldest first & Outcome\\
\midrule
0 & 0 & 0 & Empty & Budget exhausted\\
2 & 4 & 2 & $[0,3]$ & Budget exhausted\\
4 & 7 & 4 & $[0,3]$ & Budget exhausted\\
11 & 4 & 11 & $[3,2]$ & Budget exhausted\\
20 & 4 & 20 & $[5,1]$ & Budget exhausted\\
29 & 4 & 29 & $[6,0]$ & Budget exhausted\\
31 & 16 & 31 & $[6,0]$ & Budget exhausted\\
32 & 17 & 32 & $[6]$ & Budget exhausted\\
33 & 18 & 33 & $[6]$ & Halt\\
\bottomrule
\end{tabular}
\caption{Fresh invocations of the valid countdown image with sufficient capacity. Equal PCs can have different state and retirement counts. Reaching HALT after consuming the final budget unit returns halt, not budget exhaustion. These are guest transitions, not native instruction counts.}
\label{tab:revision-bytecode-prefixes}
\end{table}


\Needspace{8\baselineskip}
The native unit bundles the default countdown. For additional manual experiments,
use the fixed fixture helper from the book checkout after implementing your entry. It edits bytecode
and caller controls while preserving the native instruction text:
\Needspace{6\baselineskip}
\begin{verbatim}
python3 tools/revision_vm_programs.py countdown /tmp/vm-four.asm \
  --source /path/to/your/independent/vm.asm --budget 4
\end{verbatim}
Open the generated copy and break at \texttt{vm\_ready}. The descriptor already
contains the return state; the caller has not stored \texttt{last\_status} and
\texttt{last\_steps} yet. Expect PC seven, depth two, retired four, status zero,
fault minus one and live values zero and three. The helper refuses existing
outputs and writes a receipt naming the selected image and controls. It
supplies no interpreter implementation.

\begin{figure}[!ht]
\centering
\begin{tikzpicture}[alt={Actual native descriptor after countdown with guest budget four. First-row words are PC seven, depth two, retired four and status zero. Native outlines distinguish the PC and retired count. The second row reports fault minus one and preserves inputs.}]
\node[anchor=north west,inner sep=0,font=\small\bfseries] (title) at (0,0) {PC seven is not a count of seven instructions};
\node[anchor=north west,inner sep=0] (a) at ([yshift=-2mm]title.south west) {\MARSRevisionCrop{revision/captures/bytecode/vm-budget.png}{0.99\linewidth}{0.03125}{0.31555556}{0.58203125}{0.60888889}};
\begin{scope}[shift={(a.south west)},x={($(a.south east)-(a.south west)$)},y={($(a.north west)-(a.south west)$)}]
\draw[black,line width=.7pt] (.20,.3088) rectangle (.39,.6176);\draw[black,line width=.7pt] (.58,.3088) rectangle (.77,.6176);
\end{scope}
\end{tikzpicture}
\caption{After four successful guest instructions, PC is seven and depth is two. Solid outlines distinguish PC from retired count four. Status zero and fault minus one identify budget exhaustion; live values zero and three are checked separately. Native LaTeX supplies the title, crop and outlines.}
\label{fig:revision-bytecode-native-budget}
\end{figure}


\subsection*{Prove stack requirements before committing an operation}

For a binary operation, depth must be at least two before either operand is
loaded. The older value is left and top is right. Thus pushing seven, then
two, then SUB produces five, not minus five. Load each as a signed word, form
a wide candidate, and check signed-word range before replacing the pair.
The arithmetic reasoning from the expression lab applies, but this VM uses
its own status codes.

Division first requires both operands, then rejects a zero right operand,
then checks the quotient range. Its quotient truncates toward zero. Minimum
signed word divided by minus one is arithmetic overflow. A failure leaves the
pre-instruction live prefix and depth intact; it also leaves PC and retired
count at their pre-instruction values. Inactive cells within capacity are
unspecified workspace, so preserving the live prefix does not mean clearing
or comparing every physical stack word.

DUP needs one value and room for another. OVER needs two and room for another.
Check underflow before room. With an empty stack and capacity zero, DUP reports
underflow, not overflow. With the countdown, capacity two and budget two produce
an ordinary budget return at PC four with \([0,3]\). Budget three instead attempts
DUP and reports stack overflow: status minus six, PC four, depth two, retired
two, fault four, with that same live prefix. The failed DUP has not retired.

Use a fresh \texttt{countdown} fixture with \texttt{--capacity 2 --budget 3}
to reproduce that case. The native figures show a completed interpreter at the
same public caller boundary; your empty skeleton cannot yet produce them.
Private handler registers and frames are not required to match those figures.

\begin{figure}[!ht]
  \centering
  \begin{tikzpicture}[alt={Actual native descriptor after DUP fails for lack of room. First-row words are PC four, depth two, retired two and stack-overflow status minus six. The second row contains capacity two and fault four. Native outlines identify returned execution state and the capacity/fault pair.}]
    \node[anchor=north west,inner sep=0,font=\small\bfseries] (title) at (0,0) {The failed DUP does not retire};
    \node[anchor=north west,inner sep=0] (a) at ([yshift=-2mm]title.south west) {\MARSRevisionCrop{revision/captures/bytecode/vm-capacity.png}{0.99\linewidth}{0.03125}{0.31555556}{0.58203125}{0.60888889}};
    \begin{scope}[shift={(a.south west)},x={($(a.south east)-(a.south west)$)},y={($(a.north west)-(a.south west)$)}]
      \draw[black,line width=.7pt] (.20,.3088) rectangle (.965,.6176);\draw[black,line width=.7pt,dashed] (.20,.005) rectangle (.58,.303);
    \end{scope}
  \end{tikzpicture}
  \caption{Capacity two and budget three leave PC four, depth two and retired count two. Status minus six reports stack overflow; fault is four. The lower dashed outline includes preserved input capacity two and output fault four. Separate checks confirm the pre-failure live prefix remains zero and three. Native LaTeX supplies the title, crop and outlines.}
  \label{fig:revision-bytecode-native-capacity}
\end{figure}


For \texttt{add-overflow}, the helper generates PUSH32 2147483647, PUSH8 1,
ADD and HALT. ADD fails at PC seven after two retirements, preserving both
operands. For \texttt{zero-divisor}, it generates PUSH8 7, PUSH8 0, DIV and
HALT; DIV fails at PC four after two retirements. Neither failure is a license
to execute the later HALT or report its success.

\subsection*{Retire successful instructions and check budget before fetch}

Maintain a pending next PC until the instruction succeeds. Sequential operations
add their encoded byte size; JMP selects its target; JZ chooses target or PC+2
and pops its condition on both paths. Successful retirement installs that next
PC, increments the retired count and consumes one guest budget unit. HALT also
advances and retires once, then returns halt immediately.

Budget exhaustion is checked before the next fetch. For the two-byte image
PUSH8 7, budget one returns status zero at PC two with one retired instruction.
Budget two instead attempts another fetch and reports fetch-at-end, status minus
four, PC two, retired one. The first instruction's state remains visible.
An empty image with zero budget returns budget exhaustion; with positive budget
it reports fetch at the image end. A valid one-byte HALT with budget one returns
halt, not budget exhaustion.

\Needspace{17\baselineskip}
Do the whole-image validation even when budget is zero. The guest count measures
semantic transitions, not the MIPS work needed to establish that they are safe.
The public measurement helper runs a known countdown image at budgets zero and
64 and counts the complete \texttt{run\_vm} call:
\Needspace{6\baselineskip}
\begin{verbatim}
java -jar /path/to/cli.jar grade \
  independent/work.kts \
  /path/to/your/independent/vm.asm
\end{verbatim}

\begin{table}[!ht]
\centering
\begin{tabular}{r r r r}
\toprule
Guest budget & Retired guest instructions & Native, bitmap design & Native, byte-map design\\
\midrule
0 & 0 & 734 & 599\\
64 & 33 & 2222 & 2087\\
\bottomrule
\end{tabular}
\caption{Measured complete-function work in two qualified implementations. The byte-map implementation also groups opcode classes, so the difference cannot be assigned solely to boundary representation. Setup, assembly, session creation and exit are excluded. These are simulated MIPS instruction counts, not elapsed time or hardware cycles.}
\label{tab:revision-bytecode-work}
\end{table}


These measured designs both satisfy the same behavior. The helper reproduces
the measurement procedure on your implementation; it need not reproduce their
exact counts. Grade correctness first. A faster routine that skips validation
at zero budget is implementing a different contract. A self-branch can retire
the same guest instruction repeatedly until budget exhaustion, even though its
PC never advances to a new byte.

\subsection*{Publish the result and test fresh invocations}

Return status in \texttt{v0} and retired guest count in \texttt{v1}. Every return
also writes PC, depth, retired count, status and fault in the descriptor. This
includes validation failure and zero budget. Capacity and the full stack pointer
are preserved inputs. Unlike a decoder that promises untouched output on
failure, this entry must publish a complete failure report.

Statuses zero and one mean budget exhaustion and halt; fault is minus one in
both. Validation failures are minus one for opcode, minus two for truncation
and minus three for target; they publish zero PC, depth and retired count,
with the offending instruction start as fault. Runtime failures are minus four
for fetch at end, minus five for underflow, minus six for stack overflow,
minus seven for arithmetic overflow and minus eight for zero division. They
publish the failing PC as fault and retain the preceding live stack state.
Inspect status before interpreting a negative value as either a guest operand
or an error code.

All caller access must fit the selected code slice, the 32-byte descriptor or
the selected stack capacity. Code is read-only. Only the five output fields
and selected stack capacity may be written, including temporary writes.
Guards, descriptor inputs, other arenas and cells outside capacity are preserved.
Use full addresses when adding guest offsets or stack indices. Preserve every
bit of the public saved registers, frame and stack pointers, and return through
the incoming native return address. Submission-owned constants and private
native frames are allowed. Produce no console output.

Start construction with record classification and boundary marking, then add
target validation and complete publication of errors. Establish budget/fetch
order before implementing handlers. Make push and simple stack operations work,
then checked arithmetic and the two branches. Keep the common retirement path
separate from failures so adding one handler does not accidentally charge a
failed instruction. These are design stages rather than required private labels
or a fixed frame layout.

\Needspace{7\baselineskip}
Use the six native assessments as construction checkpoints on the same source.
The structural check retains 500 original opcode/truncation failures; the target
check retains eighteen invalid-target cases. The completed-outcome check contains
421 halts or budget returns, and the runtime check contains 308 failures. Together
these are all 1,247 original direct-call trials, with their selected arenas,
offsets, capacities, budgets and saved-state sentinels retained.

Four new public replacement sequences each make three calls in one machine.
The caller check runs twenty-three complete programs, observing initialization
at \texttt{vm\_call}, the actual VM return at \texttt{vm\_ready}, and the later
\texttt{last\_status/last\_steps} stores at \texttt{vm\_stored}. Your entry must
publish the descriptor before return; it must leave those two caller-owned words
for the caller to store. The final interval also protects data after publication.

Debug the default completed caller at those three boundaries. Before the call,
expect length eighteen and guest budget sixty-four. At return, expect
\texttt{v0}=1 and \texttt{v1}=33. Inspect PC eighteen, depth one, retired count
thirty-three, status one and fault minus one in the descriptor, and value six in
the live stack prefix. Continue to \texttt{vm\_stored} and inspect the two stored
words before the exit selector changes \texttt{v0}. The native figures above
remain observations of the supplied public interface in a completed interpreter;
your private handler layout need not match it.

From the imported package directory, the optional command-line route is:
\begin{verbatim}
java -jar /path/to/cli.jar grade \\
  independent/kernel.kts independent/vm.asm
\end{verbatim}
Repeat with \texttt{independent/caller.kts}. The extra public replacement sequences
run in the native lesson. The unfinished skeleton deliberately fails. Equivalent
decoders, boundary sets and dispatch organization are accepted; exact native
instruction counts and inactive workspace contents are not graded. Finish with
ordinary Run, which should exit silently; that transcript alone cannot establish
correct guest execution.

Finally, replace an image between calls without changing its address or length.
An all-HALT image marks every byte as a start. A replacement containing operands
must not inherit those markers. Reinitialize private boundary storage and all
logical execution state from the new invocation. Reusing the old boundary map
can make a branch into an operand appear valid; resuming old output PC can skip
new instructions. Test those possibilities after both a success and an error.
Retain your byte-boundary map, a guest-state trace and an explanation of each
return field alongside your working implementation.

